\documentclass{standalone}
\usepackage{tikz}
\usetikzlibrary{intersections, calc}
\usepackage{hsbcommon}
\begin{document}

    \def\myalpha{0.5}
    \def\mybeta{0.5}
    \def\myrho{0.5}

    \def\mycesfunktion{$
        y=\left(\alpha K^{-\rho}+\beta L^{-\rho}\right)^{-\frac{1}{rho}}
    $}
    \def\cessolvedfork{$
        K=\left(-\frac{\beta L^{-\rho}-y^{-\rho}}{\alpha}\right)^{-\frac{1}{\rho}}
    $}

\begin{tikzpicture}
    \draw[name path=abscissa,->,>=latex] (0,0) -- (7,0) node [right] {L};
    \draw[name path=ordinate,->,>=latex] (0,0) -- (0,7) node [above] {K};

    \clip (-1,-1) rectangle (7,7);

    %% Step 1: Cobb-Douglas isoquants
    \ifstep{1}{
        \foreach [count=\i] \myniveau in {2,2.4,...,5}
        {
            \draw[domain=1.2:7, name path=indiff\i] plot ({\x},{
                (\myniveau/\x^\myalpha)^(1/\mybeta)});
        }
    }

    %% Step 2: CES isoquants (overlaid in red)
    \ifstep{2}{
        \begin{scope}[color=red]
            \foreach [count=\i] \myniveau in {2,2.4,...,5}
            {
                \draw[thick, name path=ces\i, domain=1.4:7] plot ({\x},{
                    (-(\mybeta*\x^(-\myrho)-\myniveau^(-\myrho))/\myalpha)^(-1/\myrho)
                });
            }
        \end{scope}
    }

\end{tikzpicture}
\end{document}
